% !TeX root = main.tex

\chapter{Optimal Transport and Wasserstein Geometry}
\section{Motivation}
Optimal transport (OT) studies the geometry of probability distributions by quantifying the cost of transporting mass from one distribution to another.
The modern importance of OT in machine learning derives from several observations:
\begin{itemize}
    \item probability distributions can be treated as geometric objects;
    \item Wasserstein distances induce meaningful topologies on spaces of measures;
    \item gradient flows in Wasserstein space produce important PDEs;
    \item generative models and policies evolve distributions during training;
    \item reinforcement learning objectives can often be interpreted as transport problems.
\end{itemize}
The goal of this phase is to introduce:
\begin{enumerate}[label=(\roman*)]
    \item Monge and Kantorovich optimal transport,
    \item Wasserstein distances,
    \item dynamic optimal transport,
    \item Wasserstein geodesics,
    \item the geometry of $W_2$ space.
\end{enumerate}
Throughout these notes we focus primarily on the quadratic Wasserstein metric $W_2$, since it is the relevant geometry for Wasserstein gradient flows and drifting-field policies.
\section{The Monge Formulation}
\subsection{Mass transport maps}
Suppose we are given two probability measures
\[
\mu,\nu \in \cP(\R^d).
\]
Intuitively:
\begin{itemize}
    \item $\mu$ represents an initial mass distribution,
    \item $\nu$ represents a target mass distribution.
\end{itemize}
The original formulation of optimal transport was introduced by Gaspard Monge in 1781.
\begin{definition}[Transport map]
A measurable map
\[
T : \R^d \to \R^d
\]
is said to transport $\mu$ to $\nu$ if
\[
T_{\#}\mu = \nu.
\]
\end{definition}
The goal is to find the transport map minimizing transportation cost.
\begin{definition}[Monge problem]
Given a cost function
\[
c : \R^d \times \R^d \to [0,\infty),
\]
the Monge optimal transport problem is
\[
\inf_T
\int_{\R^d}
 c(x,T(x))d\mu(x)
\]
subject to
\[
T_{\#}\mu=\nu.
\]
\end{definition}
Typically one chooses
\[
c(x,y)=\|x-y\|^p.
\]
\subsection{Difficulties of the Monge problem}
The Monge formulation is highly nonlinear.
Major difficulties include:
\begin{itemize}
    \item existence of transport maps may fail;
    \item the pushforward constraint is nonlinear;
    \item \explanationlink{transport_maps_do_not_split_mass}{transport maps cannot split mass}.
\end{itemize}
\begin{example}
Suppose
\[
\mu = \delta_0,
\qquad
\nu = \tfrac12\delta_{-1}+\tfrac12\delta_1.
\]
No transport map can split the mass at $0$ into two pieces.
\end{example}
These difficulties motivate the Kantorovich relaxation.
\section{The Kantorovich Formulation}
\subsection{Transport plans}
Instead of deterministic maps, Kantorovich introduced probabilistic couplings.
\begin{definition}[Transport plan]
A transport plan between $\mu$ and $\nu$ is a probability measure
\[
\gamma \in \cP(\R^d\times\R^d)
\]
such that:
\begin{align*}
\explanationlink{first_marginal_transport_plan}{(\pi_x)_{\#}\gamma} &= \explanationlink{first_marginal_transport_plan}{\mu},
\\
(\pi_y)_{\#}\gamma &= \nu,
\end{align*}
where
\[
\pi_x(x,y)=x,
\qquad
\pi_y(x,y)=y.
\]
The set of all such couplings is denoted
\[
\Pi(\mu,\nu).
\]
\end{definition}
\begin{remark}
A transport plan specifies how much mass moves from $x$ to $y$.
Unlike Monge maps, transport plans may split mass.
\end{remark}
\subsection{Kantorovich optimal transport}
\begin{definition}[Kantorovich problem]
The Kantorovich optimal transport problem is
\[
\inf_{\gamma \in \Pi(\mu,\nu)}
\int_{\R^d\times\R^d}
 c(x,y)d\gamma(x,y).
\]
\end{definition}
This is a linear optimization problem over probability measures.
\begin{theorem}[Existence of optimal plans]
Let $c$ be \definitionlink{lower_semicontinuity}{lower semicontinuous} and bounded from below. Then for every
\[
\mu,\nu \in \cP(\R^d)
\]
there exists at least one optimal transport plan.
\end{theorem}
\section{Wasserstein Distances}
\subsection{Definition}
The most important costs are powers of Euclidean distance.
\begin{definition}[$p$-Wasserstein distance]
Let $p\ge1$. For
\[
\mu,\nu \in \cP_p(\R^d),
\]
the $p$-Wasserstein distance is
\[
W_p(\mu,\nu)
=
\left(
\inf_{\gamma\in\Pi(\mu,\nu)}
\int \|x-y\|^p d\gamma(x,y)
\right)^{1/p}.
\]
\end{definition}
The most important case for modern machine learning is:
\[
W_2.
\]
\subsection{Metric structure}
\begin{theorem}
For $p\ge1$, the Wasserstein distance defines a metric on
\[
\cP_p(\R^d).
\]
\end{theorem}
Thus probability distributions become points in a metric space.
\begin{remark}
Unlike KL divergence, Wasserstein distances compare distributions geometrically.
Even distributions with disjoint supports may have finite Wasserstein distance.
\end{remark}
\subsection{Comparison with KL divergence}
Suppose
\[
\mu = \delta_0,
\qquad
\nu = \delta_1.
\]
Then:
\[
\explanationlink{w2_between_dirac_masses}{W_2(\mu,\nu)=1},
\]
while
\[
\KL(\mu\|\nu)=\infty.
\]
This sensitivity to geometry is one reason Wasserstein distances are important in generative modeling.
\section{Optimal Maps and Brenier Theory}
The quadratic Wasserstein geometry possesses remarkable structure.
\subsection{Brenier theorem}
\begin{theorem}[Brenier]
Let
\[
\mu,\nu \in \cP_2(\R^d)
\]
and assume $\mu$ is absolutely continuous.
Then there exists a unique optimal transport map
\[
T : \R^d \to \R^d
\]
for the quadratic cost
\[
c(x,y)=\|x-y\|^2.
\]
Moreover,
\[
T = \grad \varphi
\]
for a convex function
\[
\varphi : \R^d \to \R.
\]
\end{theorem}
\begin{remark}
This theorem is foundational for Wasserstein geometry.
Optimal transport maps are \explanationlink{simple_brenier_example}{gradients of convex potentials}.
\end{remark}
\subsection{Displacement interpolation}
Given the optimal map $T$, define
\[
T_t(x)
=
(1-t)x+tT(x).
\]
Then the interpolating measures
\[
\mu_t = (T_t)_{\#}\mu
\]
define the Wasserstein geodesic between $\mu$ and $\nu$.
This interpolation is called displacement interpolation.
\section{Dynamic Optimal Transport}
The static Kantorovich formulation can be rewritten dynamically.
This viewpoint is one of the most important developments in modern OT.
\subsection{Continuity equations}
Recall the continuity equation:
\[
\partial_t \rho_t
+
\divergence(\rho_t v_t)=0.
\]
Here:
\begin{itemize}
    \item $\rho_t$ is a time-dependent density,
    \item $v_t$ is a velocity field.
\end{itemize}
The velocity field transports mass continuously.
\subsection{Benamou--Brenier formulation}
\begin{theorem}[Benamou--Brenier]
Let
\[
\mu_0,\mu_1 \in \cP_2(\R^d).
\]
Then
\[
W_2^2(\mu_0,\mu_1)
=
\inf_{(\rho_t,v_t)}
\int_0^1
\int_{\R^d}
 \|v_t(x)\|^2\rho_t(x)dxdt
\]
subject to:
\begin{align*}
\partial_t \rho_t
+
\divergence(\rho_t v_t)&=0,
\\
\rho_0 &= \mu_0,
\\
\rho_1 &= \mu_1.
\end{align*}
\end{theorem}
A \prooflink{benamou_brenier}{graduate-note proof of the Benamou--Brenier formula}
is included separately.
\subsection{Interpretation}
The quantity
\[
\int \|v_t(x)\|^2\rho_t(x)dx
\]
represents kinetic energy.
Thus:
\begin{quote}
The Wasserstein distance is the minimum kinetic energy required to transport one distribution into another.
\end{quote}
This dynamic viewpoint is central for:
\begin{itemize}
    \item Wasserstein gradient flows,
    \item fluid formulations of generative models,
    \item Schr\"odinger bridges,
    \item policy transport in RL.
\end{itemize}
\section{Wasserstein Geodesics}
\subsection{Geodesic curves}
A geodesic is a shortest path in a metric space.
\begin{definition}[Geodesic in metric space]
Let $(X,d)$ be a metric space.
A curve
\[
\gamma_t : [0,1]\to X
\]
is a geodesic if
\[
d(\gamma_s,\gamma_t)
=
|t-s|d(\gamma_0,\gamma_1)
\]
for all $s,t\in[0,1]$.
\end{definition}
\subsection{Wasserstein geodesics}
The displacement interpolation
\[
\mu_t=(T_t)_{\#}\mu
\]
constructed from the Brenier map \explanationlink{wasserstein_geodesics}{defines a geodesic} in $(\cP_2(\R^d),W_2)$.
Thus:
\begin{quote}
Optimal transport equips probability space with a Riemannian-like geometry.
\end{quote}
\subsection{Velocity representation}
Along Wasserstein geodesics, the velocity field satisfies:
\[
\partial_t \rho_t
+
\divergence(\rho_t v_t)=0.
\]
The velocity field minimizing kinetic energy defines the geodesic flow.
\section{Otto Calculus and Formal Riemannian Structure}
Felix Otto observed that $W_2$ space behaves formally like an infinite-dimensional Riemannian manifold.
\subsection{Tangent vectors}
Suppose $\rho_t$ evolves according to
\[
\partial_t \rho_t
+
\divergence(\rho_t v_t)=0.
\]
The tangent vector at $\rho$ can formally be identified with
\[
\dot \rho
=
-\divergence(\rho v).
\]
\explanationlink{minimal_velocity_gradient}{Among all velocity fields generating the same tangent vector, the minimal kinetic-energy representative is gradient-valued}:
\[
v = \grad \phi.
\]
Thus tangent vectors may be represented as:
\[
\dot \rho
=
-\divergence(\rho \grad \phi).
\]
\subsection{Formal metric tensor}
Otto's formal metric is:
\[
\langle \dot \rho_1,\dot \rho_2 \rangle_\rho
=
\int_{\R^d}
 \grad \phi_1 \cdot \grad \phi_2
 \, d\rho.
\]
This provides the formal Riemannian structure underlying Wasserstein gradient flows.
\begin{remark}
The space $\cP_2(\R^d)$ is not literally a finite-dimensional manifold.
The Otto calculus is partly formal, but extraordinarily powerful.
\end{remark}
\section{Convexity Along Geodesics}
In Euclidean optimization, convexity governs gradient descent.
In Wasserstein geometry, one instead studies convexity along geodesics.
\begin{definition}[Displacement convexity]
A functional
\[
\mathcal F : \cP_2(\R^d) \to \R
\]
is displacement convex if
\[
\mathcal F(\mu_t)
\]
is convex along Wasserstein geodesics.
\end{definition}
Examples include:
\begin{itemize}
    \item entropy functionals,
    \item internal energies,
    \item interaction energies.
\end{itemize}
Displacement convexity is fundamental for:
\begin{itemize}
    \item uniqueness of gradient flows,
    \item convergence of diffusion equations,
    \item entropy dissipation.
\end{itemize}
\section{Wasserstein Geometry in Machine Learning}
\subsection{Generative modeling}
Generative models induce distributions
\[
\rho_\theta=(f_\theta)_\#\rho_0.
\]
Training changes the generated distribution through parameter updates.
The induced trajectory:
\[
\theta_t
\mapsto
\rho_{\theta_t}
\]
may sometimes be interpreted as an approximate Wasserstein gradient flow.
\subsection{Diffusion and flow models}
Diffusion models define stochastic transport processes between distributions.
Flow matching and continuous normalizing flows define velocity fields transporting probability mass.
Schr\"odinger bridges regularize optimal transport with entropy.
\subsection{Policies as evolving distributions}
In reinforcement learning:
\[
\pi(a|s)
\]
is a conditional probability distribution over actions.
Modern policy optimization methods may be interpreted as transport in probability space:
\[
\pi_t
\rightarrow
\pi_{t+1}.
\]
This perspective motivates Wasserstein formulations of RL and drifting-field policies.
\section{Summary}
The key ideas of this phase are:
\begin{enumerate}[label=(\arabic*)]
    \item Optimal transport studies the geometry of probability distributions.
    \item The Monge formulation searches for deterministic transport maps.
    \item The Kantorovich relaxation allows probabilistic couplings.
    \item Wasserstein distances metrize spaces of probability measures.
    \item The quadratic Wasserstein metric admits a dynamic fluid formulation.
    \item Wasserstein geodesics describe optimal displacement of mass.
    \item The Otto calculus provides a formal Riemannian structure on probability space.
    \item Many learning algorithms may be interpreted as evolution equations in Wasserstein space.
\end{enumerate}
The central conceptual shift is the following:
\begin{quote}
Instead of optimizing points in Euclidean space, one optimizes probability distributions in Wasserstein space.
\end{quote}
This viewpoint underlies:
\begin{itemize}
    \item Wasserstein gradient flows,
    \item diffusion and flow models,
    \item Schr\"odinger bridges,
    \item entropy-regularized optimal transport,
    \item modern formulations of reinforcement learning.
\end{itemize}
In the next phase we introduce Wasserstein gradient flows and their connections with PDEs, variational principles, and learning dynamics.
\section*{Recommended Reading}
\begin{enumerate}
    \item C. Villani, \emph{Optimal Transport: Old and New}.
    \item G. Peyr\'e and M. Cuturi, \emph{Computational Optimal Transport}.
    \item F. Santambrogio, \emph{Optimal Transport for Applied Mathematicians}.
    \item J.-D. Benamou and Y. Brenier, ``A Computational Fluid Mechanics Solution to the Monge--Kantorovich Mass Transfer Problem.''
    \item F. Otto, ``The Geometry of Dissipative Evolution Equations.''
\end{enumerate}
